%! Periodic Phenomena — Unit 6, Lesson 6.1 — Exit Ticket (Answer Key)
\documentclass[10pt]{article}
\usepackage{precalculus-article}
\usepackage{precalculus-key}   % pulls in -boxes; do NOT also load -boxes
\newcommand{\TallMath}[1]{$\displaystyle #1\rule[-1.4em]{0pt}{3.2em}$}

\begin{document}

\pageheader{Unit 6, Lesson 6.1}{Exit Ticket --- Answer Key}
\namedateperiod

\vspace{0.10in}

\begin{notesbox}{Exit Ticket --- Key}

\textbf{1.}\quad The table below shows values of a function $r$.

\vspace{4pt}
\begin{center}
\renewcommand{\arraystretch}{1.4}
\begin{tabular}{|c|c|c|c|c|c|c|c|c|}
\hline
$x$ & 0 & 1 & 2 & 3 & 4 & 5 & 6 & 7 \\
\hline
$r(x)$ & 5 & 8 & 10 & 8 & 5 & 8 & 10 & 8 \\
\hline
\end{tabular}
\end{center}

\vspace{6pt}
(a)\quad Is this relationship periodic? Circle: \textbf{YES} \quad / \quad NO

\vspace{4pt}
(b)\quad If yes, estimate the period.

\ans{Yes. Period $= 4$. The sequence 5, 8, 10, 8 repeats starting at $x = 0$ and again at $x = 4$. The smallest positive input length before the pattern restarts is 4.}

\vspace{0.12in}
\textbf{2.}\quad Name one characteristic that repeats in every period and explain why.

\ans{Sample answer: Intervals of increase repeat in every period. Because $r(x + 4) = r(x)$ for all $x$, the output values at every point in one cycle match the output values at the corresponding point in every other cycle --- so wherever the function is increasing in one period, it is increasing at the same relative position in every period.}

\end{notesbox}

\end{document}
